\chapter{الصيغ التفاضلية ومبرهنة ستوكس}\label{ch:b3:forms}

استوعبت مبرهنةٌ واحدة من مبرهنات التحليل، على مدى قرنين، سائرَ
مبرهنات صنفها: المبرهنة الأساسية في التفاضل والتكامل،
وغرين--ريمان (المبرهَن عليها في مجلد السنة الجامعية~2)، ومبرهنة
التباعد لغاوس، ومبرهنة الدوّار لكلفن--ستوكس --- فكلٌّ منها تقول
إن تكامل مشتقةٍ ما على منطقة يساوي تكامل الكائن الأصلي على
\omterm{def:b3:forms:boundary}{الحافة}. وتجعل لغة \emph{الصيغ التفاضلية} منها عبارةً واحدة،
$\int_M\dd\omega = \int_{\partial M}\omega$، وتجعل تلك العبارة قابلة للبرهان بضربة واحدة. ويبني هذا
الفصل اللغة بأمانة --- الجبر المتعدد الخطية المتناوب، \omterm{def:b3:forms:d}{والمشتقة
الخارجية}، والسحوب العكسية، \omterm{def:b3:forms:orientation}{والتوجيه}، والمكاملة على المتنوعات
الجزئية في \cref{ch:b3:submanifolds} --- ويبرهن على مبرهنة ستوكس، ويصرف أول
الشيكات: المبرهنات التكاملية الكلاسيكية، \omterm{def:b3:forms:winding}{وعدد اللف} الذي كان
يدير \cref{ch:b3:residues} سرًّا، وفي مسألة نهاية الأسبوع، مبرهنة براور في
النقطة الصامدة. وفي كل ما يلي، تعني \emph{الناعمة} $\mathcal
C^\infty$؛ وكل
تطبيق وكل صيغة ناعمان ما لم يُذكر خلاف ذلك. ولا يكلّفنا هذا أي
عموم جدير بالاقتناء عند هذا المستوى، ويطلق الأيدي.

\section{الجبر المتعدد الخطية المتناوب}

\begin{definition}\label{def:b3:forms:alternating}
ليكن $E$ فضاءً متجهيًّا حقيقيًّا بُعده $n$. \emph{الصيغة
المتناوبة $k$-الخطية}\index{صيغة متناوبة} على $E$ هي تطبيق
$\alpha\colon E^k \to \R$، خطي في كل متغيّر، يحقق $\alpha(v_1, \dots, v_k) = 0$ كلما تساوى وسيطان. ويُكتب
فضاؤها $\Lambda^k E^*$؛ وباصطلاحٍ $\Lambda^0E^* = \R$. ويفرض التناوبُ
التخالفَ: فتبديل وسيطين يغيّر الإشارة (بنشر $\alpha(\dots, v + w, \dots, v + w, \dots) = 0$)، وعمومًا
$\alpha(v_{\sigma(1)}, \dots, v_{\sigma(k)}) =
\varepsilon(\sigma)\,\alpha(v_1, \dots, v_k)$ من أجل كل تبديلة $\sigma$.
\end{definition}

\begin{example}
على $E = \R^n$: $\Lambda^1E^* = E^*$ هو \omterm{def:b3:banach:operator}{الفضاء الثنوي}؛ والمحدد في الأساس
القانوني صيغةٌ متناوبة $n$-الخطية، وستبيّن \cref{prop:b3:forms:basis} أنه
يولّد $\Lambda^nE^*$ --- وهو السبب العميق في وحدانية المحدد إلى
غاية سلّمي. ومن أجل $k > n$، $\Lambda^kE^* = \{0\}$: إذ تكون $k$ من المتجهات
تابعة، وينشرها نشرُ إحداها على امتداد الأخريات فيلغي $\alpha$
بالتناوب.
\end{example}

\begin{definition}\label{def:b3:forms:wedge}
من أجل $\ell_1, \dots, \ell_k \in E^*$، يكون \emph{جداؤها الخارجي}\index{جداء خارجي}
الصيغةَ المتناوبة $k$-الخطية
\[
(\ell_1 \wedge \dots \wedge \ell_k)(v_1, \dots, v_k)
= \det\bigl(\ell_i(v_j)\bigr)_{1 \leq i, j \leq k} .
\]
والتناوب وتعدد الخطية هما تناوب المحدد وتعدد خطيته في أعمدته.
\end{definition}

\begin{proposition}[أساس $\Lambda^k$]\label{prop:b3:forms:basis}
ليكن $(e_1, \dots, e_n)$ أساسًا للفضاء $E$ بأساس ثنوي $(e_1^*, \dots, e_n^*)$. تشكّل الصيغ
\[
e_I^* = e_{i_1}^* \wedge \dots \wedge e_{i_k}^*,
\qquad I = \{i_1 < \dots < i_k\} \subseteq \{1, \dots, n\},
\]
أساسًا للفضاء $\Lambda^kE^*$؛ ومنه $\dim\Lambda^kE^* =
\binom nk$. وصراحةً، $\alpha =
\sum_{\abs I = k}\alpha(e_{i_1}, \dots, e_{i_k})\,e_I^*$.
\end{proposition}

\begin{proof}
\emph{التوليد.} لتكن $\alpha \in \Lambda^kE^*$ ولتكن $\beta =
\sum_I\alpha(e_I)\,e_I^*$، حيث يختصر
$\alpha(e_I)$ المقدارَ $\alpha(e_{i_1}, \dots, e_{i_k})$. والطرفان $k$-الخطية ومتناوبان،
ومنه يتوافقان بمجرد أن يتوافقا على جميع الثلاثيات $k$ $(e_{j_1}, \dots, e_{j_k})$
مع $j_1 < \dots <
j_k$ (إذ يرجع تعدد الخطية إلى ثلاثيات من متجهات الأساس،
ويرجع التناوب إلى المتزايدة تمامًا). و $e_I^*(e_{j_1},
\dots, e_{j_k}) = \det(e_{i_r}^*(e_{j_s})) = \delta_{IJ}$: فمن أجل $I = J$
تكون المصفوفة المطابقة؛ ومن أجل $I \neq J$ ينعدم سطرٌ ما. ومنه
$\beta(e_J) = \alpha(e_J)$ من أجل كل $J$: أي $\beta = \alpha$. \emph{والحرية.} إذا
كان $\sum_I c_Ie_I^* = 0$، أعطى التقييم على $(e_{j_1}, \dots, e_{j_k})$ أن $c_J = 0$.
\end{proof}

\begin{definition}\label{def:b3:forms:wedgegeneral}
يمتدّ \omterm{def:b3:forms:wedge}{الجداء الخارجي} إلى تطبيق ثنائي الخطية $\Lambda^kE^* \times \Lambda^\ell E^* \to
\Lambda^{k+\ell}E^*$، معيَّنٍ
بثنائية الخطية وبالعلاقة $(e_I^*) \wedge (e_J^*) = e_I^* \wedge e_J^*$ (بالوصل وإعادة الترتيب؛
والجداء $0$ إذا كان $I \cap J \neq
\varnothing$). وهو تجميعي، و\emph{تخالفي
التدريج}:
\[
\beta \wedge \alpha = (-1)^{k\ell}\,\alpha \wedge \beta
\qquad (\alpha \in \Lambda^k, \ \beta \in \Lambda^\ell).
\]
\end{definition}

\begin{proof}[\omnameProof]
تُفحص الخاصيتان على عناصر الأساس وتُمدَّدان بثنائية الخطية.
وأما التجميعية: فيساوي كلا تقويسَي $e_I^* \wedge e_J^* \wedge e_K^*$ إسفينَ العائلة
الموصولة من الصيغ $1$-الخطية، بصيغة المحدد في \cref{def:b3:forms:wedge}
(بنشر لابلاس بالكتل). وأما قاعدة الإشارة: فتمرير كلٍّ من عوامل
$\beta$ البالغة $\ell$ عبر عوامل $\alpha$ البالغة $k$ يكلّف
إشارةً لكل مبادلة متجاورة (أي مبادلة سطرين من المحدد)، ومنه
$(-1)^{k\ell}$ في المجموع.
\end{proof}

\begin{proposition}[السحب العكسي، الحالة الخطية]
\label{prop:b3:forms:linearpullback}
يستحث تطبيقٌ خطي $u\colon E \to F$، من أجل كل $k$، التطبيقَ الخطي $u^*\colon \Lambda^kF^* \to \Lambda^kE^*$،
$(u^*\alpha)(v_1, \dots, v_k) = \alpha(u(v_1), \dots,
u(v_k))$. وهو يحقق $u^*(\alpha \wedge \beta) = u^*\alpha
\wedge u^*\beta$ و $(u \circ w)^* = w^* \circ u^*$. وفوق ذلك:
\begin{enumerate}
  \item إذا كان $\dim E = n$ و $u\colon E \to E$، ففي المستقيم
        $\Lambda^nE^*$: $u^*\alpha = (\det u)\,\alpha$.
  \item وإذا كان $\operatorname{rk}u < k$، فإن $u^* = 0$ على $\Lambda^kF^*$.
\end{enumerate}
\end{proposition}

\begin{proof}
المتطابقات الدالّية فورية من التعاريف (وأما قاعدة الجداء
فتُفحص على أساافين الصيغ $1$-الخطية بصيغة المحدد --- $\det(\ell_i(uv_j)) =
\det((u^*\ell_i)(v_j))$ ---
ثم تُمدَّد ثنائيةَ الخطية).
(1) يرسل $u^*$ الفضاءَ الأحادي البُعد $\Lambda^nE^*$
(\cref{prop:b3:forms:basis}) في نفسه، ومنه $u^*\alpha =
c\,\alpha$ مع $c$ لا يتعلق بالصيغة
$\alpha \neq 0$؛ وبالاختبار على $\alpha = e_1^* \wedge \dots \wedge e_n^*$ و $(v_j) =
(e_j)$ نجد $c = \det(e_i^*(ue_j)) = \det u$.
(2) من أجل $v_1, \dots, v_k \in E$، تقع المتجهات $u(v_1), \dots,
u(v_k)$ في صورة $u$، وبُعدها
$< k$: فهي تابعة خطيًّا، وتنعدم \omterm{def:b3:forms:alternating}{الصيغة المتناوبة} على عائلة
تابعة (بنشر المتجهة التابعة على امتداد الأخريات).
\end{proof}

\section{الصيغ التفاضلية والمشتقة الخارجية}

\begin{definition}\label{def:b3:forms:diffform}
لتكن $U \subseteq \R^n$ مفتوحة. \emph{الصيغة التفاضلية
$k$-الخطية}\index{صيغة تفاضلية} على $U$ هي تطبيق ناعم $\omega\colon U \to \Lambda^k(\R^n)^*$؛
وفي أساس \cref{prop:b3:forms:basis} (بكتابة $\dd x_i$ بدل $e_i^*$)،
\[
\omega = \sum_{\abs I = k} a_I\,\dd x_I,
\qquad
\dd x_I = \dd x_{i_1} \wedge \dots \wedge \dd x_{i_k},
\]
بمعاملات ناعمة $a_I \in \mathcal C^\infty(U)$. وفضاؤها هو $\Omega^k(U)$؛ و $\Omega^0(U) = \mathcal
C^\infty(U)$.
والصيغة $0$-الخطية دالةٌ؛ والصيغة $1$-الخطية حقلٌ من الصيغ
الخطية (مثل التفاضل $\dd f$ لدالة)؛ والصيغة $n$-الخطية هي
$a\,\dd x_1\wedge\dots\wedge\dd
x_n$، وهي المقدار المكامَل الطبيعي في \cref{ch:b3:product}.
\end{definition}

\begin{definition}[المشتقة الخارجية]\label{def:b3:forms:d}
\emph{المشتقة الخارجية}\index{مشتقة خارجية} هي التطبيق الخطي
$\dd\colon \Omega^k(U) \to \Omega^{k+1}(U)$ المعرَّف بالعلاقة
\[
\dd\Bigl(\sum_I a_I\,\dd x_I\Bigr)
= \sum_I \dd a_I \wedge \dd x_I
= \sum_I\sum_{j=1}^n \frac{\partial a_I}{\partial
x_j}\,\dd x_j \wedge \dd x_I .
\]
وعلى الصيغ $0$-الخطية تكون التفاضلَ المعتاد.
\end{definition}

\begin{theorem}\label{thm:b3:forms:dsquare}
(a) $\dd(\omega \wedge \eta) = \dd\omega \wedge \eta +
(-1)^k\,\omega \wedge \dd\eta$ من أجل $\omega \in \Omega^k$ (أي \emph{قاعدة ليبنيتز المدرَّجة}).
(b) $\dd \circ \dd = 0$.\index{d تربيع صفر@$\dd^2 = 0$}
\end{theorem}

\begin{proof}
(a) بثنائية الخطية يكفي أن نعالج $\omega = a\,\dd
x_I$، $\eta = b\,\dd x_J$. عندئذٍ $\omega \wedge \eta =
ab\,\dd x_I \wedge \dd x_J$
و
\[
\dd(\omega\wedge\eta) = (b\,\dd a + a\,\dd b) \wedge \dd x_I
\wedge \dd x_J
= (\dd a \wedge \dd x_I) \wedge (b\,\dd x_J) +
a\,\dd b \wedge \dd x_I \wedge \dd x_J,
\]
وتمرير الصيغة $1$-الخطية $\dd b$ عبر عوامل $\dd
x_I$ البالغة $k$
يكلّف $(-1)^k$ (\cref{def:b3:forms:wedgegeneral}): فيكون الحدّ الثاني $(-1)^k\,\omega \wedge \dd\eta$.
(b) من أجل دالة: $\dd(\dd a) = \sum_{i,j}
\frac{\partial^2a}{\partial x_j\partial x_i}\,\dd x_j \wedge
\dd x_i$. والمعامل متناظر في $(i,j)$ بمبرهنة
شوارتز في المشتقات المختلطة (المبرهَن عليها في مجلد السنة
الجامعية~2؛ إذ $a$ من الصنف $\mathcal C^\infty$)، بينما $\dd x_j \wedge
\dd x_i$ متخالف:
فبمزاوجة الحدّين $(i,j)$ و $(j,i)$، يتلاشى كل شيء. وأما من أجل
$\omega = \sum
a_I\dd x_I$ عامة: $\dd\dd\omega = \sum\dd(\dd a_I \wedge \dd x_I)
= \sum(\dd\dd a_I\wedge\dd x_I - \dd a_I\wedge\dd(\dd x_I))$ حسب (a)، وينعدم الحدّان ($\dd(\dd x_I) = 0$ لأن المعامل
ثابت).
\end{proof}

\begin{definition}[السحب العكسي]\label{def:b3:forms:pullback}
ليكن $\varphi\colon U \to V$ ناعمًا (حيث $U \subseteq \R^m$ و $V \subseteq \R^n$ مفتوحتان).
\emph{السحب العكسي}\index{سحب عكسي} $\varphi^*\colon \Omega^k(V) \to \Omega^k(U)$ معرَّفٌ نقطةً نقطة
بالسحب العكسي الخطي على امتداد التفاضل: $(\varphi^*\omega)_x = (D\varphi(x))^*\,\omega_{\varphi(x)}$. وعمليًّا،
يعوّض $\varphi^*$: $\varphi^*f = f \circ
\varphi$ على الدوال، و $\varphi^*(\dd y_i) =
\dd\varphi_i = \sum_j\frac{\partial\varphi_i}{\partial
x_j}\dd x_j$، و $\varphi^*(a\,\dd y_{i_1}\wedge\dots\wedge
\dd y_{i_k}) =
(a\circ\varphi)\,\dd\varphi_{i_1}\wedge\dots\wedge
\dd\varphi_{i_k}$.
\end{definition}

\begin{theorem}\label{thm:b3:forms:pullback}
(a) $\varphi^*(\omega \wedge \eta) = \varphi^*\omega \wedge
\varphi^*\eta$ و $(\psi \circ \varphi)^* = \varphi^* \circ
\psi^*$.
(b) $\varphi^*(\dd\omega) = \dd(\varphi^*\omega)$: أي إن \omterm{def:b3:forms:d}{المشتقة الخارجية} تتبادل مع كل تعويض ناعم ---
وهي المتطابقة التي تجعلها \emph{مشتقةَ} النظرية.
(c) وإذا كان $\varphi\colon U \to V$ ناعمًا بين مفتوحتين من $\R^n$ وكانت $\omega = a\,\dd y_1 \wedge \dots \wedge \dd
y_n$،
فإن $\varphi^*\omega = (a \circ
\varphi)\,\det\bigl(D\varphi\bigr)\,\dd x_1 \wedge \dots
\wedge \dd x_n$.
\end{theorem}

\begin{proof}
(a) عبارات نقطية عن السحوب العكسية الخطية (\cref{prop:b3:forms:linearpullback})، مع
قاعدة السلسلة $D(\psi\circ\varphi)(x) = D\psi(\varphi(x))\,D\varphi(x)$.
(b) من أجل صيغة $0$-الخطية $f$: يكون $\varphi^*(\dd f) = \dd f \circ
D\varphi = \dd(f \circ \varphi)$ قاعدةَ السلسلة.
ومن أجل $\omega = a\,\dd y_I$: باستعمال (a)، $\varphi^*\omega =
(a\circ\varphi)\,\dd\varphi_{i_1}\wedge\dots\wedge
\dd\varphi_{i_k}$، ومنه بقاعدة ليبنيتز
(\cref{thm:b3:forms:dsquare}(a)) وبالعلاقة $\dd\dd\varphi_{i_r} =
0$:
\[
\dd(\varphi^*\omega) = \dd(a\circ\varphi) \wedge
\dd\varphi_{i_1}\wedge\dots\wedge\dd\varphi_{i_k}
= \varphi^*(\dd a) \wedge \varphi^*(\dd y_I)
= \varphi^*(\dd a \wedge \dd y_I) =
\varphi^*(\dd\omega).
\]
(c) نقطةً نقطة هذا هو بالضبط $u^*\alpha = (\det u)\alpha$ على الصيغ ذات الدرجة
العليا (\cref{prop:b3:forms:linearpullback}(1) مع $u = D\varphi(x)$).
\end{proof}

\begin{example}[الإحداثيات القطبية]\label{ex:b3:forms:polar}
من أجل $\varphi(r, \theta) = (r\cos\theta, r\sin\theta)$: $\varphi^*\dd x = \cos\theta\,\dd r -
r\sin\theta\,\dd\theta$، $\varphi^*\dd y = \sin\theta\,\dd r
+ r\cos\theta\,\dd\theta$، ومنه
\[
\varphi^*(\dd x \wedge \dd y) =
(\cos\theta\,\dd r - r\sin\theta\,\dd\theta) \wedge
(\sin\theta\,\dd r + r\cos\theta\,\dd\theta)
= r\,\dd r \wedge \dd\theta,
\]
فيظهر محدد ياكوبي في \cref{ex:b3:product:polar} بجبر محض --- دون أي نظرية
\omterm{def:b3:measure:measure}{قياس}. ويحوّل \cref{exo:b3:forms:9} هذه الملاحظة إلى عبارة: أنه من أجل
التكاملات الموجَّهة، تكون صيغة تغيير المتغيّرات \emph{هي} صيغة
\omterm{def:b3:forms:pullback}{السحب العكسي}.
\end{example}

\section{الصيغ المغلقة والتامة؛ المبرهنة المساعدة لپوانكاريه}

\begin{definition}\label{def:b3:forms:closedexact}
تكون $\omega \in \Omega^k(U)$ \emph{مغلقة}\index{صيغة مغلقة} إذا كان
$\dd\omega = 0$، و\emph{تامة}\index{صيغة تامة} إذا كان $\omega = \dd\eta$
من أجل $\eta \in \Omega^{k-1}(U)$ ما (أي \emph{دالة أصلية} للصيغة $\omega$).
والتامة $\Rightarrow$ المغلقة بحكم $\dd^2 = 0$؛ وأما العكس
فسؤالٌ عن \emph{شكل} $U$.
\end{definition}

\begin{example}[الصيغة الزاوية]\label{ex:b3:forms:angular}
على $U = \R^2 \setminus \{0\}$،
\[
\omega_\theta = \frac{x\,\dd y - y\,\dd x}{x^2 + y^2}
\]
مغلقةٌ (بالحساب المباشر: \cref{exo:b3:forms:4}) لكنها ليست \omterm{def:b3:forms:closedexact}{تامة}: إذ
تكاملها على امتداد الدائرة الواحدية $2\pi \neq 0$، بينما تنعدم
تكاملات الصيغ \omterm{def:b3:forms:closedexact}{التامة} على امتداد المنحنيات \omterm{def:b3:forms:closedexact}{المغلقة}
(\cref{prop:b3:forms:exactloop}). ومحليًّا، $\omega_\theta = \dd\theta$ من أجل أي تعيين ناعم $\theta$
للزاوية القطبية --- ومن هنا الاسم والعائق: إذ لا يوجد أي تعيين
كهذا على $U$ كله. وتدير هذه الصيغة الواحدة عددَ اللف
(\cref{sec:b3:forms:winding})، ومن خلاله مبرهنة البواقي في \cref{ch:b3:residues}.
\end{example}

\begin{theorem}[المبرهنة المساعدة لپوانكاريه]\label{thm:b3:forms:poincare}
لتكن $U \subseteq \R^n$ مفتوحة و\emph{نجمية الشكل} بالنسبة إلى $0$. كل
صيغة $k$-الخطية \omterm{def:b3:forms:closedexact}{مغلقة} على $U$ (حيث $k \geq 1$) تكون
\omterm{def:b3:forms:closedexact}{تامة}.\index{المبرهنة المساعدة لپوانكاريه}
\end{theorem}

\begin{proof}
نبني \emph{مؤثر تماثل} خطيًّا $h\colon
\Omega^k(U) \to \Omega^{k-1}(U)$ يحقق
\begin{equation}\label{eq:b3:forms:homotopy}
\dd(h\omega) + h(\dd\omega) = \omega
\qquad (k \geq 1);
\end{equation}
فإذا كان $\dd\omega = 0$، كان $\omega = \dd(h\omega)$ وانتهينا. ومن أجل $\omega = \sum_I a_I\,\dd x_I$
نضع
\[
h\omega = \sum_{\abs I = k}\ \sum_{r=1}^{k}(-1)^{r-1}
\Bigl(\int_0^1 t^{k-1}a_I(tx)\,\dd t\Bigr)\,
x_{i_r}\,\dd x_{i_1}\wedge\dots\wedge
\widehat{\dd x_{i_r}}\wedge\dots\wedge\dd x_{i_k}
\]
(والقبعة تحذف عاملًا؛ والتكاملات ناعمة في $x$ بالاشتقاق تحت
علامة التكامل، \cref{thm:b3:lebesgue:paramdiff}، إذ يُهيمَن على جميع المشتقات على
المتراصات). وفحصُ \eqref{eq:b3:forms:homotopy} حسابٌ يُجرى مرة واحدة في العمر،
فلنجره كاملًا. نثبّت $I$ ونأخذ $\omega =
a\,\dd x_I$ (بالخطية). أولًا،
\[
\dd(h\omega) = k\Bigl(\int_0^1t^{k-1}a(tx)\dd t\Bigr)\dd x_I
+ \sum_{r=1}^k(-1)^{r-1}\sum_{j=1}^n
\Bigl(\int_0^1t^{k}\,\partial_ja(tx)\,\dd t\Bigr)
x_{i_r}\,\dd x_j\wedge\dd x_{I\setminus i_r} :
\]
إذ تجمع الزمرة الأولى الحدودَ التي تصيب فيها $\dd$ العاملَ
$x_{i_r}$ --- ويعيد الإسفين $\dd x_{i_r}\wedge\dd
x_{I\setminus i_r}$ تركيبَ $\dd x_I$ بإشارة
$(-1)^{r-1}$ تلغي المعامل السابق، وتعطي قيم $r$ البالغة $k$
العاملَ $k$ --- بينما تجمع الزمرة الثانية الحدودَ التي تصيب
فيها $\dd$ التكاملَ (إذ تُخرج قاعدة السلسلة $t\,\partial_ja(tx)$). ثم
$\dd\omega =
\sum_j\partial_ja\,\dd x_j\wedge\dd x_I$، وبتطبيق تعريف $h$ في الدرجة $k+1$، والدليل $j$ يشغل
الخانة الأولى:
\[
h(\dd\omega) = \sum_{j=1}^n\Bigl(\int_0^1t^{k}
\partial_ja(tx)\dd t\Bigr)x_j\,\dd x_I
- \sum_{j=1}^n\sum_{r=1}^k(-1)^{r-1}
\Bigl(\int_0^1t^{k}\partial_ja(tx)\dd t\Bigr)
x_{i_r}\,\dd x_j\wedge\dd x_{I\setminus i_r} .
\]
وتتلاشى المجاميع المزدوجة في $\dd(h\omega) + h(\dd\omega)$، فيساوي من ثَمّ
\[
\Bigl(\int_0^1\bigl(k\,t^{k-1}a(tx) +
t^k{\textstyle\sum_j}x_j\,\partial_ja(tx)\bigr)\dd t\Bigr)
\dd x_I
= \Bigl(\int_0^1\frac{\dd}{\dd t}\bigl(t^ka(tx)\bigr)\dd
t\Bigr)\dd x_I
= a(x)\,\dd x_I = \omega,
\]
بالمبرهنة الأساسية في التفاضل والتكامل. وقد دخلت النجمية حيث
كان عليها أن تدخل: إذ $tx \in U$ من أجل $t \in \intcc01$، بحيث
يكون $a(tx)$ ذا معنى.
\end{proof}

\begin{remark}
من أجل $k = 1$ ومن أجل $\omega = \sum_j a_j\dd x_j$، تكون الدالة الأصلية $f(x) = \int_0^1\sum_ja_j(tx)\,x_j\,\dd t$ ---
أي التكامل الخطي للصيغة $\omega$ على امتداد القطعة $[0, x]$:
فالمبرهنة هي عبارة «حقلٌ ذو مصفوفة ياكوبية متناظرة تدرّجٌ» ذات
المتغيّرات المتعددة في السنة الجامعية~2، والآن في كل درجة.
وتبيّن الصيغة الزاوية (\cref{ex:b3:forms:angular}) أن الفرضية على $U$ ليست
زخرفية: إذ إن $\R^2\setminus\{0\}$ ليست نجمية الشكل، وهناك لا يستلزم الإغلاق
التمام. وما ينجو على \omterm{def:b3:topology:topology}{مجموعة مفتوحة} عامة تقيسه \emph{كوهومولوجيا
دي رام} $H^k(U) =
\ker\dd/\operatorname{im}\dd$ --- انظر \cref{exo:b3:forms:12} من أجل أول حساب غير بديهي.
\end{remark}

\section{التوجيه والمكاملة على المتنوعات الجزئية}

تتطلب مكاملة صيغة $k$-الخطية أرضًا موجَّهة ذات بُعد $k$. وتذكّر
من \cref{ch:b3:submanifolds} (\cref{thm:b3:submanifolds:characterizations}) أن \omterm{thm:b3:submanifolds:characterizations}{متنوعة جزئية} ذات بُعد $k$ $M \subseteq \R^n$
تكون محليًّا صورةَ وسطنة ناظمية $\gamma\colon V \to M \cap W$ (حيث $V
\subseteq \R^k$ مفتوحة
و $\gamma$ \omterm{def:b3:topology:continuity}{تماثل طوبولوجي} على صورته بتفاضل متباين).

\begin{definition}\label{def:b3:forms:orientation}
\emph{التوجيه}\index{توجيه} للمتنوعة $M$ هو اختيار، من أجل كل
$p \in M$، لأحد صنفَي توجيه أسس \omterm{def:b3:submanifolds:tangent}{الفضاء المماسّ} $T_pM$، بحيث
يكون \emph{متسقًا محليًّا}: أي توجد حول كل نقطة وسطنةٌ $\gamma$
يكون إطارها الإحداثي $(\partial_1\gamma, \dots,
\partial_k\gamma)$ موجَّهًا إيجابيًّا عند كل نقطة من
مجالها. وتسمّى هذه الوسطنات \emph{مباشرة}. وتكون $M$
\emph{قابلة للتوجيه} إذا وُجد توجيه؛ ويبيّن شريط موبيوس أن هذا
قد يخفق. وجميع المتنوعات الجزئية في هذا الفصل موجَّهة.
\end{definition}

\begin{definition}[تكامل صيغة]\label{def:b3:forms:integral}
لتكن $M$ متنوعةً جزئية موجَّهة ذات بُعد $k$ ولتكن $\omega$ صيغةً
$k$-الخطية معرَّفة على \omterm{def:b3:topology:topology}{جوار} للمتنوعة $M$، مع $\operatorname{supp}\omega \cap M$ متراصة.
(a) إذا كان $\operatorname{supp}\omega \cap M \subseteq
\gamma(V)$ من أجل وسطنة مباشرة واحدة، وضعنا
\[
\int_M\omega = \int_V\gamma^*\omega
\]
--- والطرف الأيمن تكامل لوبيغ على $V$ (\cref{ch:b3:product}) للمعامل $g$
للصيغة $\gamma^*\omega = g\,\dd u_1\wedge\dots\wedge\dd u_k$، وهو متصل ذو حامل متراص.
(b) وعمومًا، نختار عددًا منتهيًا من الوسطنات المباشرة
$\gamma_i(V_i)$ تغطّي المتراصة $\operatorname{supp}\omega \cap M$ وتجزئةً للوحدة تابعةً لها
$(\chi_i)$ (\cref{lem:b3:forms:partition})، ونضع $\int_M\omega =
\sum_i\int_M\chi_i\,\omega$، وكل حدّ محسوبٌ بواسطة
(a). ومن أجل منحنٍ ($k = 1$) موسطَن بالتطبيق $\gamma\colon\intcc ab\to\R^n$ نكتب
$\int_\gamma\omega = \int_a^b\gamma^*\omega$، دون اشتراط التباين.
\end{definition}

\begin{lemma}[الاتساق]\label{lem:b3:forms:consistency}
لا يتعلق التعريف (a) بالوسطنة المباشرة، ولا يتعلق التعريف (b)
بالتغطية ولا \omterm{lem:b3:forms:partition}{بتجزئة الوحدة}. وفوق ذلك، إذا كان $\Phi$ تماثلًا
تفاضليًّا بين جوارَي متنوعتين جزئيتين موجَّهتين مع
$\Phi(M) = M'$، يحمل إطارًا مباشرًا للمتنوعة $M$ إلى إطار مباشر
للمتنوعة $M'$ عند نقطة ما من كل مركّبة من $M$، فإن $\int_{M'}\omega =
\int_M\Phi^*\omega$.
\end{lemma}

\begin{proof}
(a) لتكن $\gamma\colon V \to M$ و $\delta\colon V' \to M$ وسطنتين مباشرتين تحتوي صورتاهما
$\operatorname{supp}\omega\cap M$. والانتقال $\tau =
\delta^{-1}\circ\gamma$ تماثل تفاضلي بين الأجزاء المفتوحة
المعنية من $V, V'$ (ونعومة الانتقالات: \cref{thm:b3:submanifolds:characterizations}، عبر الوصف
البياني المحلي)، و $\gamma = \delta\circ\tau$ هناك، ومنه
$\gamma^*\omega = \tau^*(\delta^*\omega)$
(\cref{thm:b3:forms:pullback}(a)). ونكتب $\delta^*\omega =
g\,\dd u_1\wedge\dots\wedge\dd u_k$؛ عندئذٍ (\cref{thm:b3:forms:pullback}(c))
$\tau^*(\delta^*\omega) =
(g\circ\tau)\,\det(D\tau)\,\dd u_1\wedge\dots\wedge\dd u_k$. وبكون الإطارين مباشرين، يرسل $D\tau$ أساسًا موجبًا إلى
أساس موجب: ومنه $\det D\tau > 0$، أي $\det D\tau =
\abs{\det D\tau}$، وتعطي مبرهنة
تغيير المتغيّرات (\cref{thm:b3:product:changeofvar}) أن $\int(g\circ\tau)\abs{\det D\tau} = \int g$: فيتوافق التكاملان.
\emph{وهذا كل سبب وجود \omterm{def:b3:forms:orientation}{التوجيه}}: فبلا التحكم في الإشارة، يختلف
محدد ياكوبي عن قيمته المطلقة ويصير التكامل سيّئ التعريف.
(b) إذا كانت $(\chi_i)$ و $(\tilde\chi_j)$ تجزئتين مقبولتين
(بتغطيتيهما)، فحسب (a) وبالجمعية المنتهية، $\sum_i\int\chi_i\omega =
\sum_{i,j}\int\chi_i\tilde\chi_j\omega =
\sum_j\int\tilde\chi_j\omega$، وكل حدّ
مزدوج قابل للحساب في أي من الخريطتين. وأما العبارة الأخيرة:
فإذا جالت $\gamma$ على الوسطنات المباشرة للمتنوعة $M$، جالت
$\Phi\circ\gamma$ على الوسطنات المباشرة للمتنوعة $M'$ (إذ مقارنة \omterm{def:b3:forms:orientation}{التوجيه}
ثابتة محليًّا، ومثبَّتة عند نقطة واحدة لكل مركّبة)، و $(\Phi\circ\gamma)^*\omega =
\gamma^*(\Phi^*\omega)$.
\end{proof}

\begin{lemma}[تجزئات الوحدة، الحالة المتراصة]
\label{lem:b3:forms:partition}
لتكن $K \subseteq \R^n$ متراصة ولتكن $W_1, \dots, W_m$ مجموعات مفتوحة تغطّي $K$.
توجد $\chi_1, \dots, \chi_m \in
\mathcal C^\infty(\R^n)$ تحقق $0 \leq \chi_i \leq 1$، وبحوامل $\operatorname{supp}\chi_i \subseteq W_i$ متراصة، ومع
$\sum\chi_i = 1$ على \omterm{def:b3:topology:topology}{جوار} للمجموعة $K$.\index{تجزئة الوحدة}
\end{lemma}

\begin{proof}
تقع كل $x \in K$ في $W_{i(x)}$ ما مع كرة \omterm{def:b3:forms:closedexact}{مغلقة} $\bar B(x, 2r_x) \subseteq W_{i(x)}$؛ ويستخرج
التراص $x_1, \dots, x_N$ بحيث تغطّي الكرات $B(x_s, r_{x_s})$ المجموعةَ $K$. ومن أجل
كل $s$ نأخذ نتوءًا $\theta_s \in \mathcal
C^\infty$، $0 \leq \theta_s \leq 1$، مع $\theta_s = 1$ على
$\bar B(x_s, r_{x_s})$، و $\operatorname{supp}\theta_s
\subseteq B(x_s, 2r_{x_s})$ (بتمليف الدالة المميّزة للكرة ذات نصف القطر
$\frac32r_{x_s}$، \cref{thm:b3:lp:regularization}). ونسند كل $s$ إلى دليل واحد
$i(s)$ يحقق $B(x_s, 2r_{x_s}) \subseteq W_{i(s)}$ ونضع $\Theta_i = \sum_{i(s) = i}\theta_s$. وعلى \omterm{def:b3:topology:topology}{المجموعة المفتوحة}
$\Omega_0 = \{\sum_j\Theta_j > \tfrac12\} \supseteq K$، تفي الدوال $\Theta_i/\sum_j\Theta_j$ بالغرض لكنها معرَّفة هناك فقط؛
وللتعميم، لتكن $\rho \in \mathcal
C^\infty(\R^n)$ تحقق $\rho = 0$ حيث $\sum_j\Theta_j
\geq 1$ و $\rho > 0$
حيث $\sum_j\Theta_j \leq \tfrac12$ (بتمليف قصٍّ ملائم للمقدار $1 - \sum_j\Theta_j$)، ونضع
\[
\chi_i = \frac{\Theta_i}{\rho + \sum_j\Theta_j} .
\]
والمقام $> 0$ في كل مكان ويساوي $\sum_j\Theta_j$ على $\{\sum_j\Theta_j \geq 1\}$،
وهي \omterm{def:b3:topology:topology}{جوار} مفتوح للمجموعة $K$ (إذ تقع كل نقطة من $K$ في كرة ما
$B(x_s, r_{x_s})$ حيث $\theta_s = 1$)؛ وهناك $\sum_i\chi_i = 1$. والحوامل والحدود
واضحة.
\end{proof}

\section{مبرهنة ستوكس}

\begin{definition}\label{def:b3:forms:boundary}
\emph{المتنوعة الجزئية ذات البُعد $k$ بحافة}\index{متنوعة جزئية
بحافة} $M \subseteq \R^n$ هي مجموعة مغطّاة بوسطنات ناظمية من نوعين:
\emph{الخرائط الداخلية} $\gamma\colon V \to M \cap W$ حيث $V \subseteq \R^k$ مفتوحة، و\emph{خرائط
الحافة} $\gamma\colon V \cap H^k \to M
\cap W$، حيث $H^k = \{u \in \R^k : u_k \geq 0\}$ ويمتدّ $\gamma$ ناعمًا وناظميًّا
إلى المفتوحة $V$. و\emph{الحافة} $\partial M$ هي مجموعة النقاط
المبلوغة عند $u_k = 0$؛ وهي \omterm{thm:b3:submanifolds:characterizations}{متنوعة جزئية} ذات بُعد $(k-1)$ بلا
حافة، موسطَنة بالتطبيقات $u' \mapsto \gamma(u',
0)$. ويستحث توجيهُ $M$ توجيهًا على
$\partial M$ بقاعدة \emph{الناظم الخارجي أولًا}: فعند $p \in \partial M$،
يكون أساسٌ $(w_1, \dots, w_{k-1})$ للفضاء $T_p\partial M$ موجبًا إذا وفقط إذا
كان $(\nu, w_1, \dots, w_{k-1})$ أساسًا موجبًا للفضاء $T_pM$، حيث $\nu \in T_pM \setminus T_p\partial M$ يشير خارج
$M$ (وفي خريطة حافة: $\nu =
-\partial_k\gamma$، إلى غاية إضافة مركّبات مماسّية ---
فصنف \omterm{def:b3:forms:orientation}{التوجيه} لا يراها).
\end{definition}

\begin{omfigure}
\begin{tikzpicture}[scale=1.05]
  \fill[omDef!10] (-2,0) arc (180:0:2) -- cycle;
  \draw[thick, omDef] (-2,0) arc (180:0:2);
  \draw[thick, omThm] (-2,0) -- (2,0);
  \node[omDef] at (0.1,1.35) {$M$};
  \node[omThm, below] at (1.55,-0.05) {$\partial M$};
  \draw[->, thick, omExo] (1.415,1.415) -- (2.05,2.05);
  \node[omExo, right] at (2.02,2.12) {$\nu$};
  \draw[->, thick] (1.29,1.53) arc (50:75:2);
  \draw[->, thick, omExo] (-0.6,0) -- (-0.6,-0.65);
  \node[omExo, below] at (-0.6,-0.65) {$\nu$};
  \draw[->, thick] (-0.35,0) -- (0.55,0);
  \node[below] at (0.1,-0.14) {المستحثّ};
\end{tikzpicture}

{\small قاعدة الناظم الخارجي أولًا: عند كل نقطة حافة، ضع
المتجهة الخارجية $\nu$ أولًا؛ فالأسس التي تكمّلها إلى إطار موجب
للمتنوعة $M$ توجّه $\partial M$. ومن أجل مجال مستوٍ \omterm{def:b3:forms:orientation}{بالتوجيه}
المعياري، هذه هي قاعدة عكس عقارب الساعة في غرين--ريمان.}
\end{omfigure}

\begin{lemma}[ستوكس على نصف الفضاء]
\label{lem:b3:forms:halfspace}
لتكن $\eta$ صيغةً ناعمة $(k-1)$-الخطية على $\R^k$ ذات حامل
متراص. عندئذٍ
\[
\int_{H^k}\dd\eta = \int_{\partial H^k}\eta,
\]
حيث يحمل $H^k$ التوجيهَ المعياري للفضاء $\R^k$ وتحمل $\partial H^k = \{u_k = 0\} \cong \R^{k-1}$
التوجيهَ المستحثّ، وهو $(-1)^k$ مضروبًا في \omterm{def:b3:forms:orientation}{التوجيه} المعياري
للفضاء $\R^{k-1}$.
\end{lemma}

\begin{proof}
أولًا محاسبة \omterm{def:b3:forms:orientation}{التوجيه}: الناظم الخارجي على امتداد $\partial H^k$
هو $-e_k$، و
\[
\det(-e_k, e_1, \dots, e_{k-1}) =
-\det(e_k, e_1, \dots, e_{k-1}) = -(-1)^{k-1}\det(e_1,
\dots, e_k) = (-1)^k :
\]
فيكون الإطار $(e_1, \dots, e_{k-1})$ للفضاء $\partial H^k$ موجبًا من أجل
\omterm{def:b3:forms:orientation}{التوجيه} المستحثّ إذا وفقط إذا كان $k$ زوجيًّا، ومن هنا المقارنة
المذكورة. وبالخطية نأخذ $\eta
= f\,\dd u_1\wedge\dots\wedge\widehat{\dd
u_i}\wedge\dots\wedge\dd u_k$، $f \in \mathcal
C^\infty_c(\R^k)$؛ عندئذٍ $\dd\eta =
(-1)^{i-1}\,\partial_if\,\dd u_1\wedge\dots\wedge\dd u_k$ (إذ
يكلّف نقل $\dd u_i$ إلى خانته $i - 1$ مبادلة). وثمة حالتان،
كلتاهما بتونيلي--فوبيني (\cref{thm:b3:product:fubini}) وبالمبرهنة الأساسية في
التفاضل والتكامل بمتغيّر واحد.

\emph{الحالة $i < k$.} بالمكاملة أولًا في $u_i$ على $\R$:
$\int_\R\partial_if\,\dd u_i = 0$ (بالحامل المتراص)، ومنه $\int_{H^k}\dd\eta = 0$. وتضييق $\eta$ على
$\{u_k = 0\}$ يحتوي العامل $\dd u_k$، وهو يتضيّق إلى $0$ (إذ
$u_k$ ثابت هناك): ومنه $\int_{\partial H^k}\eta =
0$ أيضًا.

\emph{الحالة $i = k$.} بالمكاملة أولًا في $u_k$ على
$\intco0\infty$:
\begin{align*}
\int_{H^k}\dd\eta &= (-1)^{k-1}\int_{\R^{k-1}}
\Bigl(\int_0^\infty\partial_kf\,\dd u_k\Bigr)\dd u' \\
&= (-1)^{k-1}\int_{\R^{k-1}}\bigl(0 - f(u', 0)\bigr)\dd u'
= (-1)^k\int_{\R^{k-1}}f(u',0)\,\dd u' .
\end{align*}
وفي جهة \omterm{def:b3:forms:boundary}{الحافة}، يتضيّق $\eta$ إلى $f(u',
0)\,\dd u_1\wedge\dots\wedge\dd u_{k-1}$، وبكون \omterm{def:b3:forms:orientation}{التوجيه}
المستحثّ $(-1)^k$ مضروبًا في المعياري، $\int_{\partial H^k}\eta =
(-1)^k\int_{\R^{k-1}}f(u',0)\,\dd u'$. فيتوافق
الطرفان.
\end{proof}

\begin{theorem}[ستوكس]\label{thm:b3:forms:stokes}
لتكن $M \subseteq \R^n$ متنوعةً جزئية متراصة موجَّهة ذات بُعد $k$ بحافة،
مع $\partial M$ حاملةً التوجيهَ المستحثّ، ولتكن $\omega$ صيغةً
ناعمة $(k-1)$-الخطية على \omterm{def:b3:topology:topology}{جوار} للمتنوعة $M$.
عندئذٍ\index{مبرهنة ستوكس}
\[
\int_M\dd\omega = \int_{\partial M}\omega .
\]
وعلى وجه الخصوص، إذا كان $\partial M = \varnothing$: $\int_M\dd\omega = 0$.
\end{theorem}

\begin{proof}
نغطّي المتراصة $M$ بعدد منتهٍ من صور الخرائط المباشرة (داخلية
أو حافّية)، ونأخذ تجزئةً للوحدة $(\chi_i)$ تابعةً للمجموعات
المفتوحة الموافقة $W_i$ من $\R^n$ (\cref{lem:b3:forms:partition})، مع $\sum\chi_i
= 1$ على
\omterm{def:b3:topology:topology}{جوار} للمتنوعة $M$. وعلى ذلك \omterm{def:b3:topology:topology}{الجوار} $\dd(\sum\chi_i) = 0$، ومنه
\[
\int_M\dd\omega = \sum_i\int_M\dd(\chi_i\omega),
\qquad
\int_{\partial M}\omega =
\sum_i\int_{\partial M}\chi_i\omega :
\]
والطرفان جمعيان، ويكفي أن نبرهن على المبرهنة من أجل صيغة حاملها
في صورة خريطة واحدة.

\emph{الخريطة الداخلية.} إذا كان $\operatorname{supp}\omega \cap M
\subseteq \gamma(V)$ مع $V$ مفتوحة في
$\R^k$: مدّدنا $\eta =
\gamma^*\omega$ بالصفر إلى $\R^k$ (فتكون ناعمة ذات حامل
متراص في $V$) وطبّقنا \cref{lem:b3:forms:halfspace} مع الحامل بعيدًا عن
$\partial H^k$ (بترحيل $V$ إلى نصف الفضاء العلوي المفتوح ---
أو بمجرد تكرار حساب الحالة $i < k$ على $\R^k$ كله): $\int_M\dd\omega =
\int_{\R^k}\dd\eta = 0$،
و $\int_{\partial M}\omega = 0$ لأن $\omega$ تنعدم بجوار $\partial M$.

\emph{خريطة \omterm{def:b3:forms:boundary}{الحافة}.} إذا كان $\operatorname{supp}\omega\cap M
\subseteq \gamma(V \cap H^k)$: فمع $\eta = \gamma^*\omega$ ممدَّدة
بالصفر،
$\gamma^*(\dd\omega) = \dd\eta$
(\cref{thm:b3:forms:pullback}(b))، ومنه حسب \cref{def:b3:forms:integral}:
\[
\int_M\dd\omega = \int_{H^k}\dd\eta
\overset{\text{\cref{lem:b3:forms:halfspace}}}{=}
\int_{\partial H^k}\eta .
\]
ويبقى أن نعيّن الطرف الأيمن بالمقدار $\int_{\partial
M}\omega$. \omterm{def:b3:forms:boundary}{والحافة} موسطَنة
بالتطبيق $\beta(u') =
\gamma(u', 0)$، و $\beta^*\omega$ هو تضييق $\eta$ على
$\{u_k = 0\}$ (أي \omterm{def:b3:forms:pullback}{السحب العكسي} بالضمّ $u'
\mapsto (u', 0)$ مركَّبًا مع
$\gamma$). ومقارنة \omterm{def:b3:forms:orientation}{التوجيه} هي $(-1)^k$ نفسها في الطرفين: إذ
يجلس الإطار $(\partial_1\beta, \dots, \partial_{k-1}\beta)$ في \omterm{def:b3:forms:orientation}{التوجيه} المستحثّ للحافة $\partial M$
بالإشارة $\det(-e_k, e_1, \dots, e_{k-1}) = (-1)^k$ المحسوبة في الخريطة (فتُسحب المتجهة الخارجية
عكسيًّا إلى $-e_k$)، وهي بالضبط الإشارة التي تربط \omterm{def:b3:forms:orientation}{التوجيه}
المستحثّ للفضاء $\partial H^k$ بالمعياري في $\R^{k-1}$
(\cref{lem:b3:forms:halfspace}). فيلغي الاصطلاحان بعضهما: $\int_{\partial H^k}\eta = \int_{\partial M}\omega$.
\end{proof}

\begin{example}[المبرهنات الكلاسيكية]
\label{ex:b3:forms:classical}
لتكن $D \subseteq \R^2$ مجالًا متراصًّا بحافة ناعمة، موجَّهًا معياريًّا.
ومن أجل $\omega = P\,\dd x +
Q\,\dd y$: $\dd\omega = \bigl(\partial_xQ -
\partial_yP\bigr)\dd x\wedge\dd y$، وتُقرأ ستوكس
\[
\int_D\Bigl(\frac{\partial Q}{\partial x} -
\frac{\partial P}{\partial y}\Bigr)\dd x\,\dd y
= \oint_{\partial D}P\,\dd x + Q\,\dd y :
\]
أي \emph{غرين--ريمان}، المبرهَن عليها من أجل المجالات الابتدائية
في مجلد السنة الجامعية~2 والآن في عمومها الطبيعي. وفي $\R^3$،
تعطي ستوكس مطبَّقةً على صيغة التدفق $2$-الخطية لحقل متجهي على
مجال متراص \emph{مبرهنة التباعد} $\int_\Omega\operatorname{div}F =
\int_{\partial\Omega}\langle F, \nu\rangle\,\dd S$، ومطبَّقةً على صيغة
$1$-الخطية على سطح بحافة، مبرهنةَ الدوّار الكلاسيكية
\emph{لكلفن--ستوكس}؛ ويفصّل \cref{exo:b3:forms:7} المعجمين معًا.
\end{example}

\section{عدد اللف}\label{sec:b3:forms:winding}

\begin{definition}\label{def:b3:forms:winding}
ليكن $\gamma\colon\intcc01\to\R^2\setminus\{a\}$ منحنيًا مغلقًا ناعمًا. \emph{عدد اللف}\index{عدد
اللف} له حول $a$ هو
\[
\operatorname{Ind}_\gamma(a) =
\frac1{2\pi}\int_\gamma\omega_\theta^a,
\qquad
\omega_\theta^a = \frac{(x - a_1)\,\dd y - (y - a_2)\,\dd
x}{(x - a_1)^2 + (y - a_2)^2} .
\]
\end{definition}

\begin{proposition}\label{prop:b3:forms:windinginteger}
$\operatorname{Ind}_\gamma(a) \in \Z$؛ وبوصفه دالةً في $a$ يكون ثابتًا على كل مركّبة مترابطة
من $\R^2\setminus\gamma(\intcc01)$ ومعدومًا على المركّبة غير المحدودة. ومن أجل $\gamma(t) = a + r(\cos2\pi t, \sin2\pi t)$:
$\operatorname{Ind}_\gamma(a) = 1$.
\end{proposition}

\begin{proof}
نأخذ $a = 0$ ونكتب $\gamma = (x, y)$، $\rho =
\norm\gamma > 0$، $\vartheta(t) =
\int_0^t\gamma^*\omega_\theta$، بحيث يكون $\vartheta' =
\frac{xy' - yx'}{x^2 + y^2}$.
وبالاصطلاح العقدي لتكن $u(t)
= \gamma(t)\,\rho(t)^{-1}\eu^{-\iu\vartheta(t)}$؛ عندئذٍ $\abs u = 1$ ويعطي حساب
مباشر
\[
\frac{u'}{u} = \frac{\gamma'}{\gamma} - \frac{\rho'}{\rho}
- \iu\vartheta'
= \frac{\bar\gamma\gamma'}{\abs\gamma^2} -
\frac{\rho'}{\rho} - \iu\,\frac{xy' -
yx'}{\rho^2}
= \frac{(xx' + yy')}{\rho^2} - \frac{\rho'}{\rho} = 0,
\]
لأن $\bar\gamma\gamma' = (xx' + yy') + \iu(xy' - yx')$ و $\rho\rho' = xx' + yy'$. ومنه تكون $u$ ثابتة: $\gamma(t) =
\rho(t)\,c\,\eu^{\iu\vartheta(t)}$ مع
$\abs c = 1$، ويفرض $\gamma(1) = \gamma(0)$ مع $\rho(1) = \rho(0)$ أن $\eu^{\iu\vartheta(1)} = \eu^{\iu\vartheta(0)} = 1$: أي $\vartheta(1) \in 2\pi\Z$،
أي $\operatorname{Ind}_\gamma(0) \in \Z$. وبوصفه دالةً في $a$ على المتممة المفتوحة للمنحني
المتراص، يكون التكامل المعرِّف \omterm{def:b3:topology:continuity}{متصلًا} (\cref{thm:b3:lebesgue:paramcont}، بالهيمنة
على \omterm{def:b3:topology:topology}{جوار} لكل $a$)؛ ودالةٌ متصلة ذات قيم صحيحة ثابتةٌ محليًّا،
ومن ثَمّ ثابتة على المركّبات. ومن أجل $\norm a$ كبير يكون
المقدار المكامَل من رتبة $O(1/\norm a)$ بانتظام في $t$، ومنه
يؤول الدليل إلى $0$ وينعدم على المركّبة غير المحدودة. وأما من
أجل الدائرة: $\gamma^*\omega_\theta^a = 2\pi\,\dd t$ مباشرةً.
\end{proof}

\begin{remark}
بالمطابقة $\C \cong \R^2$، $\frac{\dd z}{z} = \frac{x\,\dd x +
y\,\dd y}{x^2 + y^2} + \iu\,\omega_\theta$، ومنه $\operatorname{Ind}_\gamma(a) =
\frac1{2\iu\pi}\oint_\gamma\frac{\dd z}{z - a}$: وهذا هو
الدليل في \cref{ch:b3:residues}، وتعيد \cref{prop:b3:forms:windinginteger} البرهانَ على كونه
صحيحًا وثابتًا محليًّا بوسائل المتغيّر الحقيقي --- أي النصف
الطوبولوجي من مبرهنة البواقي، قائمًا الآن على ستوكس.
\end{remark}

\begin{proposition}\label{prop:b3:forms:exactloop}
إذا كانت $\omega = \dd f$ \omterm{def:b3:forms:closedexact}{تامة} على المفتوحة $U$ وكان $\gamma\colon\intcc01\to U$
منحنيًا مغلقًا، فإن $\int_\gamma\omega = 0$.
\end{proposition}

\begin{proof}
$\int_\gamma\dd f = \int_0^1(f\circ\gamma)'(t)\,\dd t =
f(\gamma(1)) - f(\gamma(0)) = 0$ --- إذ تعيّن قاعدة السلسلة $\gamma^*(\dd f)$ بالمقدار $(f\circ\gamma)'\,\dd t$.
\end{proof}

\begin{method}\label{met:b3:forms:method}
الحساب بالصيغ: (1) ميكِن --- إذ تُعاد ترتيب الأساافين بإشارات،
و $\dd$ يشتقّ المعاملات إلى $\dd
x_j$ جديدة، والسحوب العكسية
تعوّض؛ فثِق بالجبر، فهو يرمّز كل محدد ياكوبي. (2) ولتكامل صيغة
على \omterm{thm:b3:submanifolds:characterizations}{متنوعة جزئية}: وسطن مباشرةً، واسحب عكسيًّا، وكامل المعامل؛
\omterm{def:b3:forms:orientation}{والتوجيه} هو الفخّ الوحيد --- فافحص إطارًا واحدًا. (3) ولتبرهن
على متطابقة تكاملية، ابحث عن شكل ستوكس: هل المقدار المكامَل
تام؟ وهل المجال حافة؟ (4) ولتقارن تكاملين على متنوعتين جزئيتين
«متوازيتين»، طبّق ستوكس على المنطقة بينهما (أي حجة التشويه،
\cref{exo:b3:forms:10}). (5) وتكاملٌ غير معدوم \omterm{def:b3:forms:closedexact}{لصيغة مغلقة} يشهد بعائق
طوبولوجي --- فلا دالة أصلية، ولا ارتداد، ولا تمديد عديم الجذور:
وهكذا تقتل مسألة نهاية الأسبوع ارتدادات الكرة.
\end{method}

\section{تمارين}

\begin{exercise}[$\star$]\label{exo:b3:forms:1}
على $\R^3$، لتكن $\omega = x\,\dd y - z\,\dd x$ ولتكن $\eta =
\dd x + y\,\dd z$. احسب $\omega\wedge\eta$
و $\dd\omega$ و $\dd\eta$ و $\dd(\omega\wedge\eta)$، وتحقق من قاعدة ليبنيتز
المدرَّجة على هذا المثال.
\end{exercise}

\begin{exercise}[$\star$]\label{exo:b3:forms:2}
عيّن، على $\R^3$، التجسيدات الثلاثة للمؤثر $\dd$: من أجل
$f \in \Omega^0$، $\dd f \leftrightarrow \nabla f$؛ ومن أجل صيغة العمل $\omega_F = F_1\dd x + F_2\dd y + F_3\dd z$، $\dd\omega_F \leftrightarrow \operatorname{curl}F$؛
ومن أجل صيغة التدفق $\sigma_F = F_1\,\dd y\wedge\dd z + F_2\,\dd
z\wedge\dd x + F_3\,\dd x\wedge\dd y$، $\dd\sigma_F
\leftrightarrow \operatorname{div}F$. واستنتج من $\dd^2 =
0$
المتطابقتين $\operatorname{curl}\nabla f = 0$ و $\operatorname{div}\operatorname{curl}F = 0$.
\end{exercise}

\begin{exercise}[$\star\star$]\label{exo:b3:forms:3}
حدّد ما إذا كانت كل صيغة $1$-الخطية \omterm{def:b3:forms:closedexact}{مغلقة} أو \omterm{def:b3:forms:closedexact}{تامة} على مجالها،
واحسب دالةً أصلية حين توجد:
(a) $(2xy + z^2)\,\dd x + x^2\,\dd y + 2xz\,\dd z$ على $\R^3$؛
(b) $\dfrac{x\,\dd x + y\,\dd y}{x^2 + y^2}$ على $\R^2\setminus\{0\}$؛
(c) $\dfrac{-y\,\dd x + x\,\dd y}{x^2 + y^2}$ على نصف المستوي $\{x > 0\}$.
\end{exercise}

\begin{exercise}[$\star\star$]\label{exo:b3:forms:4}
(الصيغة الزاوية) تحقق من أن $\omega_\theta$ (\cref{ex:b3:forms:angular}) \omterm{def:b3:forms:closedexact}{مغلقة}؛
واحسب $\int_\gamma\omega_\theta$ من أجل $\gamma$ الدائرةَ ذات نصف القطر $r$ حول
$0$؛ واخلص إلى أن $\omega_\theta$ ليست \omterm{def:b3:forms:closedexact}{تامة} على $\R^2\setminus\{0\}$، وأن
$\R^2\setminus\{0\}$ ليست نجمية الشكل بالنسبة إلى أي من نقاطها (بطريقين:
عبر \cref{thm:b3:forms:poincare}، ومباشرةً من الهندسة).
\end{exercise}

\begin{exercise}[$\star\star$]\label{exo:b3:forms:5}
(صيغة المساحة لسطح فائق) لتكن $M \subseteq \R^n$ سطحًا فائقًا متراصًّا
موجَّهًا توجيهُه معطًى بحقل ناظمي واحدي $\nu$ (بحيث يكون $(w_1, \dots, w_{n-1})$
موجبًا إذا وفقط إذا كان $(\nu, w_1, \dots, w_{n-1})$ موجبًا في $\R^n$). برهن على أن
الصيغة $(n-1)$-الخطية $\sigma_\nu =
\sum_i(-1)^{i-1}\nu_i\,\dd x_1\wedge\dots\wedge\widehat{\dd
x_i}\wedge\dots\wedge\dd x_n$ تتضيّق على $M$ إلى \emph{صيغة
المساحة}: أي من أجل وسطنة مباشرة $\gamma$، $\gamma^*\sigma_\nu = \sqrt{\det G}\,\dd u_1\wedge\dots
\wedge\dd u_{n-1}$، حيث $G = ({}^t D\gamma)(D\gamma)$
مصفوفة غرام. \emph{(لاحظ أن $(\gamma^*\sigma_\nu)(e_1, \dots, e_{n-1}) = \det(\nu,
\partial_1\gamma, \dots, \partial_{n-1}\gamma)$ وربّع هذا المحدد.)} واحسب
$\int_{S^2}\sigma_\nu = 4\pi$.
\end{exercise}

\begin{exercise}[$\star\star$]\label{exo:b3:forms:6}
احسب $\int_{S^2}\omega$ من أجل $\omega = x\,\dd y\wedge\dd z
+ y\,\dd z\wedge\dd x + z\,\dd x\wedge\dd y$: (a) مباشرةً بالإحداثيات الكروية؛
(b) عبر ستوكس على الكرة الواحدية. واستنتج $\operatorname{vol}(B^3) = \frac{4\pi}3$ من مساحة
$S^2$، وعمّم:
$n\operatorname{vol}(B^n) = \operatorname{area}(S^{n-1})$،
وهو متسق مع \cref{thm:b3:product:ballvolume}.
\end{exercise}

\begin{exercise}[$\star\star$]\label{exo:b3:forms:7}
(المعجمان) استنبط بعناية من \cref{thm:b3:forms:stokes}: (a) مبرهنة التباعد في
$\R^3$ (بضمّ \cref{exo:b3:forms:2} و \cref{exo:b3:forms:5})؛ (b) ومبرهنة
كلفن--ستوكس $\int_S\langle\operatorname{curl}F, \nu\rangle\,\dd S =
\oint_{\partial S}\langle F, \tau\rangle\,\dd\ell$ من أجل سطح متراص موجَّه بحافة في $\R^3$.
وتحقق من توافق اصطلاحات \omterm{def:b3:forms:orientation}{التوجيه} على نصف الكرة العلوي المحدود
بخط الاستواء.
\end{exercise}

\begin{exercise}[$\star\star$]\label{exo:b3:forms:8}
(متطابقات غرين) من أجل $u, v$ ناعمتين على \omterm{def:b3:topology:topology}{جوار} لمجال متراص
$\Omega \subseteq \R^n$ بحافة ناعمة، برهن على
\[
\int_\Omega\bigl(u\,\Delta v + \langle\nabla u, \nabla
v\rangle\bigr) = \int_{\partial\Omega}u\,\partial_\nu
v\,\dd S,
\qquad
\int_\Omega\bigl(u\,\Delta v - v\,\Delta u\bigr) =
\int_{\partial\Omega}\bigl(u\,\partial_\nu v -
v\,\partial_\nu u\bigr)\dd S .
\]
واستنتج: أن دالةً توافقية على $\Omega$ تنعدم على
$\partial\Omega$ تنعدم انعدامًا تامًّا، وأن دالتين توافقيتين
لهما القيم الحدّية نفسها تتوافقان --- أي الوحدانية في مسألة
ديريكليه في \cref{ch:b3:conformal}.
\end{exercise}

\begin{exercise}[$\star\star$]\label{exo:b3:forms:9}
(تغيير المتغيّرات، الصيغة الموجَّهة) ليكن $\varphi\colon U
\to V$ تماثلًا
تفاضليًّا بين مفتوحتين من $\R^n$ مع $\det
D\varphi > 0$، ولتكن $f$ متصلة
ذات حامل متراص في $V$. برهن على أن متطابقة \omterm{def:b3:forms:pullback}{السحب العكسي} $\int_U\varphi^*(f\,\dd x_1\wedge\dots\wedge\dd x_n) =
\int_Vf\,\dd x_1\wedge\dots\wedge\dd x_n$
تكافئ مبرهنة تغيير المتغيّرات (\cref{thm:b3:product:changeofvar}) من أجل $\varphi$
كهذا، واشرح بالضبط أين ذهبت القيمة المطلقة لمحدد ياكوبي.
\end{exercise}

\begin{exercise}[$\star\star\star$]\label{exo:b3:forms:10}
(التشويه) لتكن $\omega$ صيغةً $2$-الخطية \emph{\omterm{def:b3:forms:closedexact}{مغلقة}} على
$\R^3\setminus\{0\}$ وليكن $S_r$ الكرةَ ذات نصف القطر $r$ ومركزها $0$. برهن
على أن $\int_{S_r}\omega$ لا يتعلق بنصف القطر $r > 0$ \emph{(بتطبيق ستوكس على
القشرة بين نصفَي قطر؛ وانتبه إلى توجيهَي \omterm{def:b3:forms:boundary}{الحافة})}. وطبّق ذلك
على \emph{صيغة الزاوية المجسمة}
\[
\omega = \frac{x\,\dd y\wedge\dd z + y\,\dd z\wedge\dd x +
z\,\dd x\wedge\dd y}{(x^2 + y^2 + z^2)^{3/2}} :
\]
وتحقق من أنها \omterm{def:b3:forms:closedexact}{مغلقة}، واحسب $\int_{S_r}\omega = 4\pi$، واخلص إلى أنها \omterm{def:b3:forms:closedexact}{مغلقة} وليست
\omterm{def:b3:forms:closedexact}{تامة} على $\R^3\setminus\{0\}$ --- أي شقيقة $\omega_\theta$ في البُعدين،
والمضمون الهندسي لقانون غاوس في الكهرباء الساكنة.
\end{exercise}

\begin{exercise}[$\star\star$]\label{exo:b3:forms:11}
ليكن $\gamma$ منحنيًا مغلقًا ناعمًا في $\R^2\setminus\{0\}$ مع $n = \operatorname{Ind}_\gamma(0)$. برهن
على $\int_\gamma\omega = n\int_{S^1}\omega$ من أجل \emph{كل} صيغة $1$-الخطية \omterm{def:b3:forms:closedexact}{مغلقة} $\omega$ على
$\R^2\setminus\{0\}$ \emph{(اكتب $\omega = c\,\omega_\theta
+ \dd f$ حسب \cref{exo:b3:forms:12})}. والتفسير: أنه على
المستوي المثقوب، يكون \omterm{def:b3:forms:winding}{عدد اللف} العائقَ الوحيد أمام انعدام
الأدوار.
\end{exercise}

\begin{exercise}[$\star\star\star$]\label{exo:b3:forms:12}
(أول حساب لدي رام) برهن على أن كل صيغة $1$-الخطية \omterm{def:b3:forms:closedexact}{مغلقة}
$\omega$ على $U = \R^2\setminus\{0\}$ تُكتب بكيفية وحيدة
\[
\omega = c\,\omega_\theta + \dd f,
\qquad c = \frac1{2\pi}\int_{S^1}\omega,\quad f \in
\mathcal C^\infty(U) :
\]
عرّف $f(p)$ بمكاملة $\omega - c\,\omega_\theta$ على امتداد مسار من $(1,0)$ إلى $p$
(بقطعة شعاعية ثم قوس دائري)، وبرهن على أن النتيجة لا تتعلق
بالاختيارات بالضبط لأن الدور على $S^1$ منعدم، وتحقق من $\dd
f = \omega - c\,\omega_\theta$.
واخلص: $H^1(\R^2\setminus\{0\}) \cong \R$، مولَّدة بالصيغة الزاوية.
\end{exercise}

\section{مسألة: مبرهنة براور في النقطة الصامدة}

\begin{problem}[{مسألة نهاية الأسبوع --- لا ارتداد ولا
هروب}]\label{pb:b3:forms:1}
تقول مبرهنة براور إن \emph{لكل \omterm{def:b3:topology:continuity}{تطبيق متصل} من الكرة الواحدية
\omterm{def:b3:forms:closedexact}{المغلقة} $\bar B = \bar B^n \subseteq \R^n$ إلى نفسها نقطةً صامدة} --- وهي إحدى مبرهنات
الرياضيات الكبرى، بنتائج تمتد من نظرية الألعاب (توازنات ناش)
إلى تحليل المصفوفات. وبرهانها بالصيغ التفاضلية أنظف البراهين
المعروفة: إذ تبيّن ستوكس أن الكرة السطحية ليست ارتدادًا للكرة
المصمتة، ويتبع كل شيء. وفي كل ما يلي، $S = S^{n-1} = \partial\bar B$، مع $n \geq 2$،
و
\[
\sigma = \sum_{i=1}^n(-1)^{i-1}x_i\,
\dd x_1\wedge\dots\wedge\widehat{\dd x_i}\wedge\dots\wedge
\dd x_n
\]
(والقبعة تحذف العامل).

\medskip
\noindent\textbf{الجزء الأول --- أداة \omterm{def:b3:measure:measure}{القياس}.}
\begin{enumerate}
  \item احسب $\dd\sigma$، واستنتج من ستوكس (\cref{thm:b3:forms:stokes}) أن
        $\int_S\sigma =
        n\operatorname{vol}(\bar B) > 0$، حيث تحمل $S$ توجيهَ حافة الكرة المصمتة.
  \item ومن أجل $n = 2$ و $n = 3$، عيّن تضييق $\sigma$ على $S$
        بصيغتَي طول القوس والمساحة (\cref{exo:b3:forms:5} مع
        $\nu(x) = x$) وأعد حساب $\int_S\sigma$ مباشرةً.
  \item لتكن $W \subseteq \R^N$ مفتوحة ولتكن $\varphi\colon W
        \to \R^n$ ناعمة مع $\norm{\varphi(x)} = 1$ من
        أجل كل $x \in W$. برهن على أن $\varphi^*(\dd\sigma) = 0$. \emph{(اشتقّ
        $\norm\varphi^2 = 1$: فتقع صورة $D\varphi(x)$ في الفضاء الفائق
        $\varphi(x)^\perp$، وبُعده $n - 1$؛ ثم طبّق \cref{prop:b3:forms:linearpullback}(2).)}
  \item وأين الخلل في «البرهان» التالي على أن $\int_S\sigma = 0$: «إن $S$
        متراصة بلا حافة، وتضييق $\sigma$ على $S$ صيغةٌ ذات
        درجة عليا عليها، ومن ثَمّ \omterm{def:b3:forms:closedexact}{مغلقة}، ومن ثَمّ $\int_S\sigma =
        \int_S\dd(\text{شيء ما}) = 0$
        بستوكس»؟ (عيّن الكلمة الخاطئة.)
  \item اشرح في فقرة واحدة استراتيجية الجزء الثاني: ماذا
        سيُكامَل، وعلى ماذا، ومن أين سيأتي التناقض.
\end{enumerate}

\medskip
\noindent\textbf{الجزء الثاني --- لا ارتداد ناعم.} لنفترض،
للتناقض، أن $r$ \emph{ارتدادٌ ناعم} للكرة المصمتة على كرتها
السطحية: أي إن $r$ ناعم على \omterm{def:b3:topology:topology}{جوار} للمجموعة $\bar B$، و $r(\bar B) \subseteq S$،
و $r(x) = x$ من أجل كل $x \in S$.
\begin{enumerate}[resume]
  \item برّر $\int_Sr^*\sigma = \int_S\sigma$ \emph{(إذ يتضيّق $r$ على $S$ إلى المطابقة:
        فإذا كان $\gamma$ وسطنةً مباشرة لقطعة من $S$، فإن
        $r\circ\gamma = \gamma$)}.
  \item باستعمال ستوكس على $\bar B$، برهن على $\int_Sr^*\sigma =
        \int_{\bar B}\dd(r^*\sigma)$.
  \item برهن على $\dd(r^*\sigma) = r^*(\dd\sigma) = 0$ (بالسؤال 3 مطبَّقًا على
        $\varphi = r$)، واخلص: \emph{أنه لا يوجد أي ارتداد
        ناعم $\bar B \to S$.}
  \item وسوِّ الحالة المستبعدة $n = 1$ يدويًّا: برهن مباشرةً
        على أنه لا يوجد أي \omterm{def:b3:topology:continuity}{تطبيق متصل} $\intcc{-1}1 \to
        \{-1, 1\}$ يثبّت الطرفين
        معًا، وسمِّ المبرهنة التي استعملتها.
\end{enumerate}

\medskip
\noindent\textbf{الجزء الثالث --- براور الناعمة.} لتكن $g$
ناعمة على \omterm{def:b3:topology:topology}{جوار} للمجموعة $\bar B$ مع $g(\bar B)
\subseteq \bar B$ و\emph{لا} نقطة
صامدة في $\bar B$.
\begin{enumerate}[resume]
  \item برهن على $\delta = \min_{x\in\bar B}\norm{g(x) - x} >
        0$.
  \item من أجل $x \in \bar B$ لتكن $u(x) = \frac{x -
        g(x)}{\norm{x - g(x)}}$ ولتكن $r(x) = x +
        t(x)\,u(x)$ تقاطعَ
        نصف المستقيم $\{x + tu(x) : t \geq 0\}$ مع $S$. حُلّ المعادلة من الدرجة
        الثانية $\norm{x + tu}^2 = 1$ واحصل على
        \[
        t(x) = -\langle x, u(x)\rangle +
        \sqrt{1 - \norm x^2 + \langle x, u(x)\rangle^2}
        \;\geq\; 0 .
        \]
  \item برهن على أن المقدار تحت الجذر موجب \emph{تمامًا} على
        $\bar B$: فإذا كان $1 - \norm x^2 + \langle x,
        u(x)\rangle^2 = 0$، فإن $\norm x = 1$ و $\langle
        x, u(x)\rangle = 0$،
        أي $\langle x, x -
        g(x)\rangle = 0$، أي $\langle x, g(x)\rangle =
        1$؛ وبكوشي--شوارتز مع $\norm x = 1$،
        $\norm{g(x)} \leq 1$، وهذا يفرض $g(x) = x$ --- وهو مستبعد. واستنتج
        أن $r$ ناعم على \omterm{def:b3:topology:topology}{جوار} للمجموعة $\bar B$.
  \item برهن على $r(\bar B) \subseteq S$ وعلى $r(x) = x$ من أجل $x
        \in S$
        \emph{(من أجل $\norm x = 1$، تحقق من $t(x) = 0$
        باستعمال $\langle x, u(x)\rangle \geq 0$، وهو نفسه ينتج من $\langle x, x - g(x)\rangle = 1 -
        \langle x, g(x)\rangle \geq 0$)}. واخلص مع
        الجزء الثاني: \emph{أن لكل تطبيق ذاتي ناعم للكرة
        المصمتة $\bar B$ نقطةً صامدة.}
\end{enumerate}

\medskip
\noindent\textbf{الجزء الرابع --- براور المتصلة.} لتكن $f\colon\bar B\to\bar B$
متصلة بلا نقطة صامدة.
\begin{enumerate}[resume]
  \item برهن على $\varepsilon = \min_{\bar B}\norm{f - x} > 0$، وأنتج تطبيقًا \emph{حدوديًّا} $p\colon\R^n\to
        \R^n$
        يحقق $\sup_{\bar B}\norm{p - f} <
        \varepsilon/2$ (بستون--فايرشتراس، \cref{thm:b3:complete:stoneweierstrass}، إحداثيًّا
        إحداثيًّا --- وبرّر الانتقال من التقريب السلّمي إلى
        المتجهي).
  \item وقد يغادر التطبيق $p$ الكرةَ؛ فنضع $g = \frac{p}{1
        + \varepsilon/2}$. برهن على
        $g(\bar B) \subseteq \bar B$ وعلى $\sup_{\bar B}\norm{g - f} < \varepsilon$.
  \item استنتج تناقضًا مع الجزء الثالث واخلص: \emph{أن لكل \omterm{def:b3:topology:continuity}{تطبيق
        متصل} $\bar B^n \to \bar B^n$ نقطةً صامدة.}
  \item بيّن بمثال أن المبرهنة تخفق على: الكرة المفتوحة؛
        والكرة السطحية $S$؛ والتاج المغلق. وأي خاصية للمجموعة
        $\bar B$ يفقدها كل مثال مضاد؟
\end{enumerate}

\medskip
\noindent\textbf{الجزء الخامس --- العوائد.}
\begin{enumerate}[resume]
  \item (بيرون--فروبنيوس، الوجود) لتكن $A$ مصفوفةً من الرتبة
        $n\times n$ جميع عناصرها $> 0$، ولتكن $\Delta = \{x \in \R^n : x_i \geq 0,\ \sum x_i =
        1\}$. برهن على
        أن التطبيق $x \mapsto Ax/\norm{Ax}_1$ معرَّف جيدًا ومتصل على $\Delta$، وأن
        $\Delta$ مطابقة طوبولوجيًّا لكرة \omterm{def:b3:forms:closedexact}{مغلقة} من $\R^{n-1}$
        (\omterm{def:b3:topology:continuity}{بتماثل طوبولوجي} شعاعي من محدّبة متراصة ذات داخل غير
        خالٍ في غلافها التآلفي)، واخلص إلى أن للمصفوفة $A$
        متجهةً ذاتية بعناصر موجبة تمامًا وقيمةً ذاتية $> 0$.
  \item استنتج أن لكل مصفوفة عشوائية بعناصر موجبة (أي بأعمدة
        مجموعها $1$) متجهةً احتمالية ساكنة $\pi = A\pi$ ---
        أي المتجهة من نمط بيج رانك. (وتصح الوحدانية أيضًا لكنها
        تحتاج إلى أدوات أخرى.)
  \item (الكرة المشعرة، الإعداد) لتكن $v$ ناعمة على \omterm{def:b3:topology:topology}{جوار}
        للمجموعة $S = S^{n-1}$ مع $\langle v(x),
        x\rangle = 0$ و $\norm{v(x)} = 1$ من أجل
        $x \in S$ (أي \emph{حقل مماسّ واحدي}). ومن أجل
        $t \in \R$ نضع $F_t(x) = x + t\,v(x)$. برهن على $\norm{F_t(x)} = \sqrt{1 + t^2}$ على $S$: أي إن
        $F_t$ يرسل $S$ في الكرة السطحية $\sqrt{1+t^2}\,S$.
  \item برهن على أن $P(t) = \int_SF_t^*\sigma$ \emph{كثير حدود} في $t$ \emph{(إذ
        كل دالة معامل للصيغة $F_t^*\sigma$ في خريطة كثيرُ حدود
        في $t$، بمعاملات ناعمة في متغيّر الخريطة؛ والمكاملة
        خطية)}.
  \item برهن على أنه من أجل $\abs t$ صغير، يكون $F_t$ تماثلًا
        تفاضليًّا من $S$ على $\sqrt{1+t^2}\,S$: بالتباين من أجل $t\operatorname{Lip}(v) < 1$؛
        وبالتماثل التفاضلي المحلي بمبرهنة الدالة العكسية
        (\cref{thm:b3:submanifolds:ift} في الخرائط)؛ وبكون الصورة مفتوحة \omterm{def:b3:forms:closedexact}{ومغلقة} في
        الكرة السطحية الهدف المترابطة. واستنتج، باستعمال
        \cref{lem:b3:forms:consistency} والتحجيم $\sigma_{\lambda x} =
        \lambda^{n}\,\sigma_x$ تحت $x \mapsto \lambda x$ (تحقق منه)، أن
        \[
        P(t) = \pm(1 + t^2)^{n/2}\int_S\sigma,
        \qquad\text{بالإشارة} + \text{ من أجل } t
        \]
        (\omterm{def:b3:forms:orientation}{والتوجيه} محفوظ بالاتصال انطلاقًا من $t = 0$).
  \item اخلص (ميلنور): أنه إذا كان $n$ \emph{فرديًّا}، فإن
        $(1 +
        t^2)^{n/2}$ ليس كثير حدود في $t$، ومع ذلك يوافق كثير الحدود
        $P(t)/\int_S\sigma$ بجوار $0$ --- وهو تناقض. ومن ثَمّ فإن \emph{الكرات
        السطحية ذات البُعد الزوجي $S^{n-1}$ (حيث $n$ فردي) لا
        تحمل أي حقل مماسّ واحدي}، ولا تحمل، بالتعيير والتنعيم
        (بالالتفاف إحداثيًّا ثم بالإسقاط --- وبرّر الخطوتين)،
        أي حقل مماسّ متصل عديم الجذور البتة: فكل ريح على الأرض
        تترك نقطةً ساكنة.

  \item (الكرات الفردية تُمشَّط بحرية) أبرز على $S^{2m-1}
        \subseteq \R^{2m} \cong \C^m$ حقلًا
        مماسًّا واحديًّا ناعمًا صريحًا: $v(x) = \iu x$
        بالاصطلاح العقدي، أي
        \[
        v(x_1, y_1, \dots, x_m, y_m) =
        (-y_1, x_1, \dots, -y_m, x_m) .
        \]
        وتحقق من التماس ومن الطول الواحدي، واخلص إلى أن ثنائية
        التماثل في السؤال 23 مضبوطة: أي إن الكرة السطحية قابلة
        للتمشيط إذا وفقط إذا كان بُعدها فرديًّا. وأين ينكسر
        برهانُ كثير الحدود في السؤال 22 من أجل $n$ زوجي؟
  \item (الجذور من السلوك على \omterm{def:b3:forms:boundary}{الحافة}) لتكن $f \colon \bar
        B^n \to \R^n$ متصلة مع
        $\langle f(x),
        x\rangle \geq 0$ من أجل كل $x \in S^{n-1}$. برهن على أن $f$
        تنعدم في موضع ما من $\bar B^n$. \emph{(وإلا فإن $g(x) = -f(x)/\norm{f(x)}$
        يرسل $\bar B$ اتصاليًّا في $S \subseteq \bar B$؛ فطبّق براور على $g$
        وناقض فرضية \omterm{def:b3:forms:boundary}{الحافة}.)} واستنتج \emph{محك الغمر}: أن
        تطبيقًا \omterm{def:b3:topology:continuity}{متصلًا} $F\colon\R^n\to\R^n$ يحقق $\frac{\langle F(x), x\rangle}{\norm x} \to
        +\infty$ حين $\norm x \to \infty$ يكون
        غامرًا --- وهو السلف المنتهي البُعد لحجج القسرية في
        التحليل غير الخطي.
\end{enumerate}
\end{problem}
